\section{Discussion}
	
	The present study sought to determine whether individual differences in Orthographic Sensitivity and Language Proficiency dissociably modulate the timing and nature of morphological processing, using ERP and behavioral measures to evaluate predictions that group-level analyses have been unable to adjudicate. Three findings from the results bear directly on this question and structure the discussion that follows. First, individual differences dissociated from behavioral responses but were clearly visible in neural dynamics, a finding that underscores the importance of ERP measures for characterising variation in morphological processing that is invisible at the level of response times. Second, while the predicted clean dissociation between Orthographic Sensitivity and Language Proficiency was only partially supported, both dimensions modulating early nonword processing at the N250, Language Proficiency emerged as the selective driver of the three-way Complexity $\times$ Family Size $\times$ LP interaction at the N400, providing the clearest evidence that depth of lexical-semantic knowledge shapes a qualitatively distinct aspect of morphological processing. Third, the N250 and N400 showed qualitatively different patterns of skill modulation—the complexity effect reversed at high skill at the N250 but showed monotonic attenuation without reversal at the N400, suggesting that early parsing efficiency and later semantic integration are governed by distinct mechanisms that respond differently to accumulated lexical experience. We consider each of these findings in turn.
	
	\subsection{Individual Differences Dissociate from Behavioral Responses but Are Visible in Neural Dynamics}
	
	Behavioral response times replicated well-established morphological frequency effects. For words, both base frequency and family size facilitated lexical decision latencies, with effect sizes consistent with the extensive prior literature on morphological processing in visual word recognition. For nonwords, morphological complexity produced a robust decision cost—complex nonwords were rejected more slowly than simple nonwords—and base frequency produced a surprising inhibitory effect, such that nonwords built on higher-frequency stems were responded to more slowly. This base frequency inhibition is consistent with the view that high-frequency stems engage stronger lexical expectations that are violated when no whole-word entry is retrieved, increasing decision uncertainty and prolonging response times.
	
	Against this backdrop of clear lexical and morphological effects, individual differences in Orthographic Sensitivity and Language Proficiency had limited influence on response latencies. OS showed a marginal and unstable trend toward faster responding for words, and a reliable but non-morphological main effect for nonwords; LP did not reliably predict response times in either condition. Crucially, neither dimension modulated the structure of morphological effects in behavior. This null pattern at the behavioral level replicates findings in the individual differences literature \citep{CITATION} showing that readers with different orthographic and semantic profiles do not differ reliably in overall lexical decision speed but do differ systematically in the \emph{pattern} of effects sensitive to morphological structure.
	
	The ERP measures resolved what the behavioral measures could not. Both N250 and N400 amplitudes were systematically modulated by individual differences in ways that were component-specific, morphological-variable-specific, and—for the N400 in nonwords—three-way in their structure. The dissociation between insensitive behavior and informative neural dynamics is theoretically important: it indicates that the mechanisms by which individual differences express themselves operate within the processing cascade—at sequential stages of form encoding and meaning retrieval—rather than at the point of response preparation. Response times collapse across this cascade, producing a single latency value that averages over processing dynamics that are only recoverable in the time-resolved ERP signal. This finding speaks directly to the methodological claim that behavioral and electrophysiological measures tap qualitatively different aspects of reading, and argues for the routine inclusion of neural measures in studies of individual differences in morphological processing.
	
	\subsection{N250 Effects and the Partial Dissociation at the Form-Processing Stage}
	
	The most striking finding at the N250 was not the presence of a complexity effect—which is predicted by all current accounts—but its reversal at high levels of both Orthographic Sensitivity and Language Proficiency. For readers scoring one standard deviation above the mean on either dimension, complex nonwords elicited less negative N250 amplitudes than simple nonwords, indicating that the presence of a real morphological suffix facilitated rather than impeded early processing. This pattern requires explanation, and its interpretation has consequences for how the N250 complexity effect should be understood more generally.
	
	For real words, the results additionally provided support for the role of Orthographic Sensitivity as a modulator of early form-based processing. The Base Frequency $\times$ OS interaction was significant, indicating that orthographically sensitive readers show a more pronounced early neural response to high-frequency stems relative to low-frequency stems compared with readers lower in OS. The Family Size $\times$ OS interaction showed a parallel marginal trend. These findings are consistent with the Lexical Quality Hypothesis \citep{CITATION} and with the morpho-orthographic account \citep{CITATION}: readers with more precise sublexical orthographic representations encode stem identity more efficiently, producing stronger activation signatures at the stage indexed by the N250—prior to full lexical access. The direction of these interactions, with OS amplifying rather than reducing morphological sensitivity, suggests that orthographic precision enables more complete and differentiated engagement with morphological form cues, consistent with a system in which the quality of form representations determines the resolution with which morphological structure can be registered early in the recognition cascade.
	
	The dissociation between OS and LP was considerably less clean for nonwords. Both dimensions modulated the N250 complexity effect in a nearly parallel fashion—complexity costs at the N250 were large at low levels of both OS and LP, attenuated at mean levels, and reversed at high levels. The reversal itself is theoretically noteworthy: for highly skilled readers, the presence of a legitimate derivational suffix relative to a pseudosuffix reduced rather than increased early processing costs. This is inconsistent with any account that treats suffix recognition as an additional parsing operation that always incurs a processing cost. Instead, it suggests that for skilled readers the form-cue value of a real suffix—its predictable orthographic structure and its association with morphological relationships learned across the lifetime of reading experience—provides positive scaffolding that facilitates early form processing. This interpretation aligns with the discriminative learning framework \citep{CITATION}, which holds that form-meaning mappings are learned through accumulated error-driven weight adjustment; under this view, the strength of the association between a suffix form and its processing consequences varies with the richness of prior exposure, and the advantage for real suffixes at high skill would reflect precisely this learned facilitation.
	
	The convergence of OS and LP at the nonword N250 does not negate the theoretical distinction between orthographic and semantic processing dimensions—it rather suggests that both dimensions contribute complementary resources to the same early mechanism when processing must generalize to novel forms. Without a stored whole-word representation to retrieve, the processing system must rely simultaneously on orthographic familiarity with the suffix form and on the breadth of the learned morphological vocabulary to bootstrap the early parse. In this context, both dimensions of reader skill serve a common function—they enrich the form-level input representations that drive the early morpho-orthographic response—even though they are orthogonal in the behavioral domain.
	
	\subsection{N400 Effects and Language Proficiency as the Selective Modulator of Semantic Integration}
	
	At the N400, the predicted role of Language Proficiency as the selective modulator of semantic processing received its clearest support in the nonword analyses. The three-way Complexity $\times$ Family Size $\times$ LP interaction was significant, indicating that LP modulated the joint influence of morphological complexity and morphological family size on N400 amplitude. Specifically, at low LP the complexity effect was substantially amplified for large-family relative to small-family nonwords, reflecting disproportionately strong semantic processing demands when both morphological complexity and family richness are simultaneously high. This interaction attenuated progressively with increasing LP and was absent at high LP, where complexity effects were reduced and family size no longer differentiated complex from simple nonwords during semantic integration. This pattern indicates that for highly proficient readers, the semantic consequences of morphological complexity are resolved efficiently regardless of family size—that is, LP confers the capacity to draw on rich morphological family representations without incurring elevated processing cost. The selective role of LP in this interaction, with OS entering only as a non-significant main effect covariate, provides the clearest empirical support for the predicted dissociation: Language Proficiency specifically governs how morphological family information is recruited during semantic integration.
	
	For words, the N400 analyses revealed a significant Base Frequency $\times$ LP interaction, indicating that more proficient readers derive greater semantic benefit from high-frequency stems during lexical access. This aligns with the theoretically motivated prediction and extends the nonword result to familiar lexical items: depth of lexical-semantic knowledge amplifies the processing advantage associated with more frequently encountered morphological forms. By contrast, the Family Size $\times$ LP interaction for words did not reach significance. This asymmetry is interpretable within a multi-stage account: for familiar words with stored whole-word representations, semantic integration of morphological family information proceeds robustly regardless of LP level—the family size effect reflects a property of the stored lexical entry rather than an inference that must be drawn online. In contrast, for nonwords lacking whole-word entries, LP becomes the critical determinant of whether and how family size information is recruited during semantic integration, because the system must generalize to a novel form in real time.
	
	Across both word models at the N400, Orthographic Sensitivity yielded a significant main effect independently of morphological structure—higher OS was associated with more positive N400 responses overall. This finding, which was not predicted by the theoretical framing, suggests that orthographic precision affects the general efficiency of semantic access rather than the specific recruitment of morphological information. One interpretation is that more precise orthographic representations reduce the processing load associated with form encoding at earlier stages, freeing resources for lexical-semantic integration and producing globally smaller (more positive) N400 responses as a consequence. This is consistent with the cascade view of visual word recognition \citep{CITATION}, in which efficiency gains at the orthographic processing stage propagate forward to facilitate subsequent semantic access.
	
	\subsection{Theoretical Implications, Limitations, and Conclusion}
	
	The present results speak to the ongoing debate between decomposition-based accounts and discriminative learning frameworks of morphological processing. The finding that morphological sensitivity at both the N250 and N400 varies continuously as a function of orthographic and semantic experience is difficult to reconcile with accounts positing a fixed, obligatory decomposition procedure: a mechanism of this kind would not be expected to produce reader-dependent differences in early parsing efficiency, nor would it predict the reversal of complexity effects at high skill levels observed in the nonword N250. The reversal is particularly telling in this regard—a pure parsing account predicts greater, not lesser, early engagement for structurally complex stimuli, whereas a distributional learning account predicts exactly this pattern: real suffix sequences become processing facilitators because they have been encountered repeatedly in predictable morphological relationships, and the strength of that facilitation scales with accumulated exposure. This finding is, to our knowledge, among the clearest electrophysiological demonstrations that what presents as a morphological effect can be understood as a consequence of form-cue learning rather than structural parsing.
	
	At the same time, the temporal structure of the word-level effects—base frequency modulating the N250 and family size modulating the N400, with OS selectively amplifying the former and LP the latter—is consistent with the multi-stage architecture proposed by morpho-orthographic accounts \citep{CITATION}. The present results do not cleanly adjudicate between a classical parsing mechanism and a learned statistical sensitivity, but they demonstrate that whichever mechanism is responsible, its expression is shaped by accumulated lexical experience in a temporally structured and component-specific way. The combination of findings points toward a reading system that is neither purely combinatorial nor purely form-driven, but one in which learned orthographic and semantic representations interact across time to produce the observed pattern of morphological sensitivity.
	
	Several limitations warrant consideration. The ERP analyses relied on median-split categorization of base frequency and family size to yield stable averaged waveforms, which introduces some loss of statistical power relative to the continuous regression framework used for reaction times; future work using single-trial ERP regression methods \citep{CITATION} would allow the full range of variation in both morphological predictors and individual differences to be exploited simultaneously. Real word stimuli were not fully counterbalanced across morphological variables, as base frequency, family size, and morphological complexity covary in naturalistic distributions, and the PCA dimensions—though orthogonal—are composites of behavioral measures that are imperfect proxies for the underlying cognitive constructs they are interpreted as indexing. Finally, the sample consisted of college-age proficient readers of English, limiting generalization to developing readers or populations with atypical reading profiles for whom the orthographic-semantic balance may be differently configured.
	
	The present study provides electrophysiological evidence that individual differences in reading skill are expressed in the temporal dynamics of morphological processing in ways that are not captured by overt response times. Language Proficiency emerged as the selective driver of semantic integration at the N400, particularly for novel morphologically complex forms where no stored whole-word representation is available. Orthographic Sensitivity modulated early form-based processing at the N250 for real words, and both dimensions jointly modulated early processing for novel forms, suggesting that early morpho-orthographic parsing is sensitive to accumulated lexical experience broadly rather than form-based knowledge exclusively. Together, these findings underscore the importance of time-resolved neural measures for characterising morphological processing, demonstrate that the orthographic-semantic balance established as an individual difference dimension in the behavioral literature has a genuine and temporally structured neural correlate, and point toward a reading system in which the representation and use of morphological information is shaped continuously by accumulated experience with print.